% Options for packages loaded elsewhere
\PassOptionsToPackage{unicode}{hyperref}
\PassOptionsToPackage{hyphens}{url}
\documentclass[
]{article}
\usepackage{xcolor}
\usepackage[margin=1in]{geometry}
\usepackage{amsmath,amssymb}
\setcounter{secnumdepth}{-\maxdimen} % remove section numbering
\usepackage{iftex}
\ifPDFTeX
  \usepackage[T1]{fontenc}
  \usepackage[utf8]{inputenc}
  \usepackage{textcomp} % provide euro and other symbols
\else % if luatex or xetex
  \usepackage{unicode-math} % this also loads fontspec
  \defaultfontfeatures{Scale=MatchLowercase}
  \defaultfontfeatures[\rmfamily]{Ligatures=TeX,Scale=1}
\fi
\usepackage{lmodern}
\ifPDFTeX\else
  % xetex/luatex font selection
\fi
% Use upquote if available, for straight quotes in verbatim environments
\IfFileExists{upquote.sty}{\usepackage{upquote}}{}
\IfFileExists{microtype.sty}{% use microtype if available
  \usepackage[]{microtype}
  \UseMicrotypeSet[protrusion]{basicmath} % disable protrusion for tt fonts
}{}
\makeatletter
\@ifundefined{KOMAClassName}{% if non-KOMA class
  \IfFileExists{parskip.sty}{%
    \usepackage{parskip}
  }{% else
    \setlength{\parindent}{0pt}
    \setlength{\parskip}{6pt plus 2pt minus 1pt}}
}{% if KOMA class
  \KOMAoptions{parskip=half}}
\makeatother
\usepackage{graphicx}
\makeatletter
\newsavebox\pandoc@box
\newcommand*\pandocbounded[1]{% scales image to fit in text height/width
  \sbox\pandoc@box{#1}%
  \Gscale@div\@tempa{\textheight}{\dimexpr\ht\pandoc@box+\dp\pandoc@box\relax}%
  \Gscale@div\@tempb{\linewidth}{\wd\pandoc@box}%
  \ifdim\@tempb\p@<\@tempa\p@\let\@tempa\@tempb\fi% select the smaller of both
  \ifdim\@tempa\p@<\p@\scalebox{\@tempa}{\usebox\pandoc@box}%
  \else\usebox{\pandoc@box}%
  \fi%
}
% Set default figure placement to htbp
\def\fps@figure{htbp}
\makeatother
\setlength{\emergencystretch}{3em} % prevent overfull lines
\providecommand{\tightlist}{%
  \setlength{\itemsep}{0pt}\setlength{\parskip}{0pt}}
\usepackage{bookmark}
\IfFileExists{xurl.sty}{\usepackage{xurl}}{} % add URL line breaks if available
\urlstyle{same}
\hypersetup{
  hidelinks,
  pdfcreator={LaTeX via pandoc}}

\author{}
\date{\vspace{-2.5em}}

\begin{document}

\section{Jop Rubén Escobedo Cruz}\label{jop-rubuxe9n-escobedo-cruz}

\section{Matrícula 1597754}\label{matruxedcula-1597754}

\section{Fecha de entrega ----
17/09/2026}\label{fecha-de-entrega---17092026}

\section{Objetivos de la práctica}\label{objetivos-de-la-pruxe1ctica}

\section{1--- Inicair una sesión en
Rstudio}\label{inicair-una-sesiuxf3n-en-rstudio}

\section{2--- Realziar una descripción estadística de un conjunto de
datos.}\label{realziar-una-descripciuxf3n-estaduxedstica-de-un-conjunto-de-datos.}

\section{3--- Comprender el diseño del panel de
RStudio.}\label{comprender-el-diseuxf1o-del-panel-de-rstudio.}

\section{4--- Conozca la sintaxis
markdown.}\label{conozca-la-sintaxis-markdown.}

\section{Parte 1: Importar datos ----}\label{parte-1-importar-datos--}

\section{Importar desde archivos csv
----}\label{importar-desde-archivos-csv--}

\section{Función read.csv importa datos de excel a
R}\label{funciuxf3n-read.csv-importa-datos-de-excel-a-r}

getwd()

trees \textless- read.csv(``DBH\_1.csv'', header=TRUE) head(trees)

\section{Ingresar datos directo en la consola
----}\label{ingresar-datos-directo-en-la-consola--}

dbh \textless- c(16.5, 25.3, 22.1, 17.2, 16.1, 8.1, 34.3, 5.4, 5.7,
11.2, 24.1, 14.5, 7.7, 15.6, 15.9, 10, 17.5, 20.5, 7.8, 27.3, 9.7, 6.5,
23.4, 8.2, 28.5, 10.4, 11.5, 14.3, 17.2, 16.8)

\section{Datos de URL no seguras (http)
----}\label{datos-de-url-no-seguras-http--}

url \textless-
``\url{https://www.agricultura.gob.mx/sites/default/files/sagarpa/document/2025/07/11/2245/14072025-actividades-inspeccion-movilizacion.csv}''
inspeccion \textless- read.csv(url) head(inspeccion)

prof\_url\_2 \textless-
paste0(``\url{https://www.agricultura.gob.mx/sites/default/files/sagarpa/}'',
``document/2025/07/11/2245/'',
``14072025-actividades-inspeccion-movilizacion.csv'')

senasica \textless- read.csv(prof\_url\_2) head(senasica)

\section{Datos de URL seguras (https): Dropbox
----}\label{datos-de-url-seguras-https-dropbox--}

library(repmis) \# No olvidar cargar la paquetería

conjunto \textless-
source\_data(``\url{https://www.dropbox.com/s/hmsf07bbayxv6m3/cuadro1.csv?dl=1}'')
head(conjunto) \# muestra las primeras seis filas de la BD

\section{Datos de URL seguras (https): Github
----}\label{datos-de-url-seguras-https-github--}

library(readr)

file \textless-
paste0(``\url{https://raw.githubusercontent.com/mgtagle/}'',
``202\_Analisis\_Estadistico\_2020/master/cuadro1.csv'')

inventario \textless- read\_csv(file) head(inventario)

\section{Parte 2: Operaciones con la base de datos
----}\label{parte-2-operaciones-con-la-base-de-datos--}

mean(trees\$dbh) \# El signo de \$ informa que necesitamos la columna
dbh sd(trees\$dbh)

\section{Selección mediante restricciones
----}\label{selecciuxf3n-mediante-restricciones--}

\section{Indica la sumatoria de los individuos en el objeto tree con un
dbh \textless{} a
10}\label{indica-la-sumatoria-de-los-individuos-en-el-objeto-tree-con-un-dbh-a-10}

sum(trees\$dbh \textless{} 10)

which(trees\$dbh \textless{} 10)

\section{Excluir los diámetros de la parcela 2. El símbolo ! indica
NO}\label{excluir-los-diuxe1metros-de-la-parcela-2.-el-suxedmbolo-indica-no}

trees.13 \textless- trees{[}!(trees\$parcela==``2''),{]} trees.13

\section{Selección de una submuestra
----}\label{selecciuxf3n-de-una-submuestra--}

trees.1 \textless- subset(trees, dbh \textless= 10) head(trees.1)

mean(trees\(dbh)
mean(trees.1\)dbh)

\section{Histogramas de la muestra y la submuestra (Figura 5)
----}\label{histogramas-de-la-muestra-y-la-submuestra-figura-5--}

\section{Argumentos de hist y etiquetas vistos en clase
(Revisar\_conjunto\_datos.R)}\label{argumentos-de-hist-y-etiquetas-vistos-en-clase-revisar_conjunto_datos.r}

hist(trees\$dbh, main = ``Muestra original trees'', xlab = ``DBH'', ylab
= ``Frecuencia'')

hist(trees.1\$dbh, main = ``dbh \textless{} 10 cm. trees.1'', xlab =
``DBH'', ylab = ``Frecuencia'')

\section{Parte 3: Representación gráfica
----}\label{parte-3-representaciuxf3n-gruxe1fica--}

\section{Histogramas ----}\label{histogramas--}

mamiferos \textless-
read.csv(``\url{https://www.openintro.org/data/csv/mammals.csv}'')

hist(mamiferos\$total\_sleep)

hist(mamiferos\$total\_sleep, \# Datos xlim = c(0,20), ylim = c(0,14),
\# Cambiar los límites de x \& y main = ``Total de horas sueño de las 39
especies'', \# Cambiar el título xlab = ``Horas sueño'', \# Cambiar eje
de las x ylab = ``Frecuencia'', \# Cambiar eje de las y las = 1, \#
Cambiar orientación de y col = ``navajowhite'') \# Cambiar color de las
barras

\section{Barplot o gráfico de barras
----}\label{barplot-o-gruxe1fico-de-barras--}

data(``chickwts'') head(chickwts{[}c(1:2,42:43, 62:64), {]})

feeds \textless- table(chickwts\$feed) \# Tabla de frecuencias por tipo
de alimento feeds

barplot(feeds) barplot(feeds{[}order(feeds, decreasing = TRUE){]})

\section{Figura 10: barplot horizontal
----}\label{figura-10-barplot-horizontal--}

\section{Ejercicio del laboratorio: ``Ahora te toca a
tí''}\label{ejercicio-del-laboratorio-ahora-te-toca-a-tuxed}

barplot(feeds{[}order(feeds){]}, \# Orden creciente: la mayor queda
arriba horiz = TRUE, \# Barras horizontales las = 1, \# Etiquetas
horizontales main = ``Frecuencias por tipos de\n alimentación'', \#
Título en dos líneas xlab = ``Número de Pollos'', \# Eje de las x col =
c(``\#FFCC00'', ``\#FFD11A'', ``\#FFD633'', ``\#FFDB4D'', ``\#FFE066'',
``\#FFE680'')) \# Gradiente amarillo

\end{document}
